
\begin{unnumbered}

\section{ВСТУП}

У сучасному суспільстві батьки та опікуни щодня стикаються з необхідністю систематичного відстеження медико-біологічних показників новонародженої дитини: режиму годування та сну, динаміки маси тіла й зросту, графіка вакцинацій і вікових досягнень. \gls{ВООЗ} наголошує, що регулярний моніторинг фізичного розвитку протягом перших двох років життя є ключовим інструментом раннього виявлення відхилень і своєчасного медичного втручання. Попри широкий вибір мобільних застосунків для моніторингу немовляти, жоден із них не поєднує в собі повноцінного антропометричного моніторингу відповідно до стандартів \gls{ВООЗ}, синхронізації між пристроями в реальному часі. Це зумовлює актуальність розроблення комплексної кросплатформної системи, яка усуне зазначені недоліки.

Метою дипломного проєкту є розроблення кросплатформної клієнт-серверної системи агрегації та обробки медико-біологічних даних з використанням хмарних технологій --- мобільного застосунку для відстеження розвитку немовляти на платформах iOS та Android. Об'єктом дослідження є процеси агрегації, зберігання та аналізу медико-біологічних даних про розвиток немовляти. Предметом дослідження є методи та інструменти розроблення кросплатформних мобільних застосунків із застосуванням хмарних технологій.

У процесі виконання роботи буде проведено аналіз предметної області та огляд існуючих аналогів, сформульовано вимоги до системи. Також буде здійснено огляд та обґрунтування вибору технологічного стеку: React Native та Expo для клієнтської частини, Supabase з PostgreSQL для хмарного бекенду. Наступні етапи передбачають проєктування реляційної схеми бази даних із захистом на рівні рядків (\gls{RLS}), реалізацію всіх модулів застосунку, а також тестування системи на рівнях юніт-тестів та функціонального тестування.

Розроблений застосунок є актуальним практичним рішенням для широкого кола користувачів --- від молодих батьків, які вперше стикаються з необхідністю ведення медичних записів дитини, до сімей, де обидва батьки потребують спільного доступу до актуальних даних у режимі реального часу. Застосунок розгорнуто на хмарній інфраструктурі та готовий до використання на обох платформах.

\end{unnumbered}
